%SSS - additional limitation is assumption of perfect repair and predictable and unchanging repair time 
%L: read by SSS - keeping for future reference - I also think that having a predictable and unchanging repair time could be too much of a limitation to the model. I can relax on this limitation by changing the repair time to a random variable generated from a probability distribution function. This probability distribution function can be found using the historical completion time of successful repairs
%SSS - This can introduce a lot of error to your recovery comparison, because we may underestimate the repair time and end up without resources we had assumed are available, or overestimate it and end up with poor utilization. The latter is a smaller problem. 
%L: some recovery process models published in recent years must know each component's exact repair time (For example, sarkale's). However, predicting the exact repair time for every component in different failure events is challenging. Therefore, using the probability distribution of the repair time learned from previous failure events is the most promising approach to finding the most possible outcome of the recovery process with the Monte Carlo methods (that I am already using). The benefit of using this method is the repair time of the component will be randomly sampled from the probability distribution function (rounded up to the nearest integer value in the unit of the time step) in every simulation of the simulation. In some simulations, we may underestimate the repair time; in others, we may overestimate. When we collected a large number of simulations (1,000 in my case study), we could find the most possible recovery rate of the system based on the average system performance at each time step of every simulation simulated. So we can tell if one repair scheme is better than the other or having how many repair crews is enough. Furthermore, we can find the lower and upper bounds of recovery time from the simulations that completed the earliest and latest. 


%SSS One idea is to assume a fixed ceiling for each particular repair, so your estimate of repair will be conservative, possibly low utilization, but at least feasible.


%SSS - I still have an issue with looking at repair crews as equal and interchangeable elements